\scanpage{149}
% archive foliation: 148
Compléments aux notes IHES

\deleted{Exercice}
\uline{Remarque} La méthode de démonstration de 1.7. prouve plus
\unsure{généralt} l'énoncé suivant : Soit $X$ et $Y$ sur $S$,
$X$ propre \inserted{surjectif} séparable à fibres
$\mathrm{uni}^t$ connexes, et $Y$ connexe,
\inserted{les préschémas loc. noeth.}
On choisit un pt géométrique $c$ de $X\times_S Y$, d'où des pts géom
$a,b,d$ dans $X,Y,S$ et un hom canonique
\[
\pi_1(X\times_S Y,c)\longrightarrow
\pi_1(X,a)\mathbin{\times_{\pi_1(S,d)}}\pi_1(Y,b).
\]
Cet hom est surjectif. \deleted{Si}
Il est bijectif si
\inserted{$Y$ a une section sur $S$, plus généralement si}
la condition suivante est vérifiée : Il existe un $S'$
\deleted{\unsure{étale}}\inserted{connexe}
non vide, \deleted{\unsure{étale}} sur $S$, tel que
$Y\times_S S'=Y'$ ait une section sur $S'$
[i.e. \deleted{\unsure{il faut}} il existe un $S$-morphisme $S'\to Y$]
\deleted{et que \unsure{les fibres} \unsure{aient une composante connexe}.}
et que $\pi_1(X')\mathbin{\unsure{-|\longrightarrow}}\pi_1(X)$ soit
injectif [\unsure{par exemple si} $S'$ étale sur $S$]
\deleted{\illegible{1}}.
Cet énoncé s'applique en particulier chaque fois que $S$ est le spectre
d'un corps $k$, et donne
\[
\pi_1(X\times_k Y)\xrightarrow{\sim}
\pi_1(X)\mathbin{\times_{\pi_1(k)}}\pi_1(Y)
\]
[$X$ propre $\mathrm{uni}^t$ connexe sur $k$, $Y$
\deleted{\unsure{uni}} connexe sur $k$ sans plus]

\margin{\unsure{$S'\to S$}}
\margin{\deleted{\unsure{$X'\to X$}}}
